% PAGES: 293-294
% Continues the sentence left open at the end of sec09_c3_1.tex (page 292 /
% folio 276 ends mid-word: "...invested in the"). No environment was open
% across that boundary, so this file simply resumes the same sentence.
second asset, or invested in the minimum variance portfolio fully invested in
the two assets.

(iii) What is the percentage difference between the VaRs obtained by using
the formulas (9.93) and (9.92)?

\begin{answerx}
(i) Let $\VaR_1(10\text{ days},98\%)$, $\VaR_2(10\text{ days},98\%)$,
$\VaR_m(10\text{ days},98\%)$ be the 10--day 98\% VaR of the \$100 million
portfolio fully invested in the first asset, fully invested in the second
asset, and invested in the minimum variance portfolio fully invested in the
two assets, respectively.

Let $\sigma_1 = 0.25$, $\sigma_2 = 0.35$, and $\rho_{1,2} = -0.25$ be the
standard deviations and the correlation of the returns of the two assets,
respectively. From (9.82), we find that the standard deviation of the return
of the minimum variance portfolio fully invested in the two assets is
$\sigma_m = 0.177139$.

From the approximate VaR formula (9.93), we obtain that
\[
\begin{aligned}
\VaR_1(10\text{ days},98\%) &\approx \sqrt{\frac{10}{252}}\,\sigma_1 z_{98}\cdot\$100\text{mil} \;=\; \$10.228\text{mil};\\
\VaR_2(10\text{ days},98\%) &\approx \sqrt{\frac{10}{252}}\,\sigma_2 z_{98}\cdot\$100\text{mil} \;=\; \$14.319\text{mil};\\
\VaR_m(10\text{ days},98\%) &\approx \sqrt{\frac{10}{252}}\,\sigma_m z_{98}\cdot\$100\text{mil} \;=\; \$7.247\text{mil};
\end{aligned}
\]
recall that $z_{98}=2.0537$.

The VaR of the minimum variance portfolio is significantly smaller that the
VaR of both the portfolio invested in the first asset, and of the portfolio
invested in the second asset, which shows that portfolio diversification can
be used to reduce the risk of the portfolio.
% NOTE: "significantly smaller that the VaR" -- source prints "that", not
% "than" (folio 277); reproduced as printed.

(ii) Let $\mu_1 = 0.03$ and $\mu_2 = 0.06$ be the expected values of the
returns of the two assets, respectively.

Since $z_{98} = 2.0537$, we obtain from (9.92) that
\[
\begin{aligned}
\VaR_1(10\text{ days},98\%) &= \sqrt{\frac{10}{252}}\,\sigma_1 z_{98}\cdot\$100\text{mil} \;-\; \frac{10}{252}\mu_1\,\$100\text{mil} \;=\; \$10.109\text{mil};\\
\VaR_2(10\text{ days},98\%) &= \sqrt{\frac{10}{252}}\,\sigma_2 z_{98}\cdot\$100\text{mil} \;-\; \frac{10}{252}\mu_2\,\$100\text{mil} \;=\; \$14.081\text{mil}.
\end{aligned}
\]

The expected value of the return of the minimum variance portfolio fully
invested in the two assets is $\mu_m = w_1\mu_1 + w_2\mu_1$,
% NOTE: source prints $w_2\mu_1$ here (folio 277); the second term should
% read $w_2\mu_2$ on definitional grounds ($\mu_m$ is the weighted average
% $w_1\mu_1+w_2\mu_2$), but this is reproduced exactly as printed and not
% corrected, per the fidelity rule.
where $w_1$ and $w_2$ are given by (9.80--9.81). We obtain that $\mu_m =
0.04107$, and, from (9.92), it follows that
\[
\VaR_m(10\text{ days},98\%) \;=\; \sqrt{\frac{10}{252}}\,\sigma_m z_{98}\,\$100\text{mil} \;-\; \frac{10}{252}\mu_m\,\$100\text{mil} \;=\; \$7.084\text{mil}.
\]

(iii) The percentage differences between the VaRs obtained by using the
formulas (9.93) and (9.92) are all smaller than 2.5\%, as follows:

for the portfolio fully invested in the first asset:
\[
\frac{\lvert 10.228-10.109\rvert}{10.228} \;=\; 0.0116 \;=\; 1.16\%.
\]

for the portfolio fully invested in the second asset:
\[
\frac{\lvert 14.319-14.081\rvert}{14.319} \;=\; 0.0166 \;=\; 1.66\%.
\]

for the minimum variance portfolio fully invested in the two assets:
\[
\frac{\lvert 7.247-7.084\rvert}{7.247} \;=\; 0.0225 \;=\; 2.25\%.
\]

In other words, the VaRs obtained using the expected values of the returns of
the two assets and the VaRs obtained using only the standard deviations of
the returns are very similar. Thus, in order to estimate Value at Risk, there
is no need to find the expected values of the returns of assets or
portfolios, which is very difficult to do with any accuracy in practice. This
is similar to option pricing in the Black--Scholes framework, where the drift
of the asset (i.e., the expected value of the instantaneous return of the
asset) is not needed in order to compute the Black--Scholes values of the
options.
\hfill$\square$
\end{answerx}

The usefulness of Value at Risk comes from its simplicity: it is a single
number that makes comparing VaRs at different times easy, and makes
deteriorating trends easy to spot. However, its simplicity is also its main
drawback, since VaR does not distinguish between possible sizes of the losses
beyond the $1-C$ percentile.

A coherent risk measure that distinguishes between different sizes of losses
under tail events is expected shortfall, also called conditional Value at
Risk (cVaR). The $N$ days expected shortfall of a portfolio at confidence
level $C$ is the expected loss of the portfolio over the next $N$ days
conditional on the loss being greater than $\VaR(N,C)$. Expected shortfall is
subadditive, unlike VaR; see section 9.6.1 for details on the subadditivity
of VaR.

\subsection{VaR of combined portfolios and subadditivity}

An intuitive requirement for a risk measure is that it would decrease with
portfolio diversification. Formally, this is called the subadditivity
requirement for a coherent risk measure: If $\mathrm{rm}(V)$ denotes the risk
measure of portfolio $V$, then this risk measure is subadditive if and only
if $\mathrm{rm}(V_1 + V_2) \le \mathrm{rm}(V_1) + \mathrm{rm}(V_2)$ for any two
portfolios $V_1$ and $V_2$.

Value at Risk is not a coherent risk measure, as first pointed out by Artzner
et al. \cite{6}; see the example below. This has potentially very detrimental
consequences for using VaR in practice. For example, it may not be true that
the overall Value at Risk of a financial institution is smaller than the sum
of the VaRs for different groups within the firm. Note that expected
shortfall (cVar) is subadditive, and it has the further advantage of better
estimating tail risk exposure.
% NOTE: this same page (folio 278) prints "conditional Value at Risk (cVaR)"
% a few lines above and "expected shortfall (cVar)" here -- the
% capitalization of the final letter is inconsistent in the source; both
% spellings are reproduced exactly as printed.

The example below of an instance when VaR is not subadditive is adapted from
Artzner et al. \cite{6}.

Assume that the changes $V_1(N)-V_1(0)$ and $V_2(N)-V_2(0)$ in the values of
two given portfolios over an $N$ day time horizon are independent and
identically distributed, with probability density functions given by
\begin{equation}
f_1(x) \;=\; f_2(x) \;=\;
\begin{cases}
0.05, & \text{if} \ \ -2 \le x \le 0;\\
0.90, & \text{if} \ \ 0 < x \le 1;\\
0, & \text{if} \ \ 1 < x.
\end{cases}
\end{equation}
